\documentclass[11pt]{article}

\usepackage[margin=1in]{geometry}
\usepackage{amsmath}
\usepackage{amssymb}
\usepackage{booktabs}
\usepackage{hyperref}

\title{Transformer Scaling Law}
\author{}
\date{}

\begin{document}
\maketitle

\section{Dense Transformer Scaling}

We fit a Chinchilla-style loss model to the dense transformer sweep:
\[
  L(N, D) = E + \frac{A}{N^\alpha} + \frac{B}{D^\beta},
\]
where \(N\) is the parameter count, \(D\) is the number of training tokens seen,
and compute is approximated as
\[
  C = 6ND.
\]

Using the 10-epoch constant-learning-rate sweep, the fitted dense law was:
\[
  L(N, D)
  \approx
  3.86 \times 10^{-14}
  + \frac{7.10 \times 10^{12}}{N^{2.761}}
  + \frac{5.04 \times 10^{6}}{D^{1.302}}.
\]

This implies the compute-optimal allocation:
\[
  N_{\mathrm{opt}}(C)
  \approx
  39.23 \left(\frac{C}{6}\right)^{0.320},
\]
\[
  D_{\mathrm{opt}}(C)
  \approx
  0.0255 \left(\frac{C}{6}\right)^{0.680}.
\]

\begin{table}[h]
\centering
\begin{tabular}{lll}
\toprule
Quantity & Value & Interpretation \\
\midrule
\(E\) & \(3.86 \times 10^{-14}\) & Near-zero irreducible loss floor \\
\(\alpha\) & \(2.761\) & Model-size loss term falls quickly with \(N\) \\
\(\beta\) & \(1.302\) & Data/token loss term falls substantially with \(D\) \\
\(N\) exponent & \(0.320\) & Parameters grow slowly with compute \\
\(D\) exponent & \(0.680\) & Training tokens grow faster with compute \\
\bottomrule
\end{tabular}
\caption{Dense transformer fitted law and optimal allocation summary.}
\end{table}

The main qualitative result is that the optimal allocation is data-heavy.
For a \(10\times\) increase in compute, the fitted law recommends roughly
\[
  10^{0.320} \approx 2.1\times
\]
more parameters, but
\[
  10^{0.680} \approx 4.8\times
\]
more training tokens. Thus, for this addition task, the dense model benefits
from scaling both model size and data, but the fitted optimum allocates most
new compute to more token exposure.

This differs from the classic language-model Chinchilla result, where model
size and data size scale more evenly. A plausible explanation is that this is
a small synthetic arithmetic task with a near-zero achievable loss. Once the
model has enough capacity to represent the algorithm, additional training
exposure can be more valuable than additional parameters. The result is also
consistent with the observed frontier: small models dominate at low compute,
then the frontier moves to larger models only once the token budget is large.

\section{MoE Scaling}

For the mixture-of-experts model, we distinguish between total parameters and
active parameters:
\[
  N_{\mathrm{total}} = \text{all stored parameters},
\]
\[
  N_{\mathrm{active}} = \text{parameters active for a token}.
\]
The sweep used top-1 routing with four experts, so the MoE models have roughly
three times as many stored parameters as active parameters:
\[
  \frac{N_{\mathrm{total}}}{N_{\mathrm{active}}} \approx 3.
\]
The active compute proxy is
\[
  C_{\mathrm{active}} = 6N_{\mathrm{active}}D.
\]

The fitted active-parameter MoE law was:
\[
  L(N_{\mathrm{active}}, D)
  \approx
  4.64 \times 10^{-18}
  + \frac{2.42 \times 10^{18}}{N_{\mathrm{active}}^{4.000}}
  + \frac{4.77 \times 10^{6}}{D^{1.284}}.
\]

This gives:
\[
  N_{\mathrm{active,opt}}(C_{\mathrm{active}})
  \approx
  203.61 \left(\frac{C_{\mathrm{active}}}{6}\right)^{0.243},
\]
\[
  D_{\mathrm{opt}}(C_{\mathrm{active}})
  \approx
  0.00491 \left(\frac{C_{\mathrm{active}}}{6}\right)^{0.757}.
\]

We also fit a stored-capacity view:
\[
  L(N_{\mathrm{total}}, D)
  \approx
  1.95 \times 10^{-13}
  + \frac{1.65 \times 10^{20}}{N_{\mathrm{total}}^{4.000}}
  + \frac{4.77 \times 10^{6}}{D^{1.284}}.
\]
The corresponding proxy optimum was:
\[
  N_{\mathrm{total,opt}}(C_{\mathrm{proxy}})
  \approx
  453.03 \left(\frac{C_{\mathrm{proxy}}}{6}\right)^{0.243},
\]
\[
  D_{\mathrm{opt}}(C_{\mathrm{proxy}})
  \approx
  0.00221 \left(\frac{C_{\mathrm{proxy}}}{6}\right)^{0.757}.
\]

\begin{table}[h]
\centering
\begin{tabular}{llll}
\toprule
Law & \(N\) exponent & \(D\) exponent & Fit note \\
\midrule
Dense & \(0.320\) & \(0.680\) & Data-heavy, but both terms identifiable \\
MoE active & \(0.243\) & \(0.757\) & Even more data-heavy \\
MoE total & \(0.243\) & \(0.757\) & Similar because \(N_{\mathrm{total}}\) tracks \(N_{\mathrm{active}}\) \\
\bottomrule
\end{tabular}
\caption{Compute-optimal exponents for dense and MoE fits.}
\end{table}

The MoE result is rougher than the dense result. The most important warning is
that \(\alpha\) hit the upper bound of the fitting range:
\[
  \alpha = 4.000.
\]
This should not be interpreted as evidence that the true MoE model-size
exponent is exactly \(4\). Rather, it means the available data does not cleanly
identify the model-size term. The optimizer wanted the \(N\)-dependent loss
penalty to vanish even faster than the allowed range. In practice, this says
that the observed MoE loss variation is explained more clearly by \(D\) than by
model size.

The empirical results support this caution. At the 10-epoch slice, dense models
still win the best validation loss at each sample size. However, MoE models do
show useful behavior at matched active-compute points. In particular, the extra
stored expert capacity helps some smaller active models:
\[
  \text{200k samples, config 0: dense } 0.531 \quad \text{vs. MoE } 0.157,
\]
\[
  \text{200k samples, config 1: dense } 0.00683 \quad \text{vs. MoE } 0.00472.
\]
This is the expected qualitative advantage of MoE: more stored capacity can help
without increasing active compute by the same factor.

It's worth noting that the current MoE uses hard top-1 routing, no
load-balancing loss, no expert-capacity control, and no auxiliary router loss.
Those choices make optimization noisier and can hide clean scaling behavior.
The task also saturates near zero loss, which makes log-space power-law fitting
unstable at the high-compute end. For a cleaner MoE scaling law, the next steps
would be to add load balancing, sweep more token budgets at high \(N\), and fit
frontier-adjacent points separately from obviously under-optimized points.

\section{Pallas Kernel}

The MoE feed-forward layer uses the sequence
\[
  h = \mathrm{ragged\_dot}(x, W_{\mathrm{up}}),\qquad
  h = \mathrm{gelu}(h),\qquad
  y = \mathrm{ragged\_dot}(h, W_{\mathrm{down}}).
\]
This materializes the hidden activation \(h\in\mathbb{R}^{M\times F}\), where
\(M\) is the number of routed tokens and \(F\) is the expert hidden width. We
implemented a Pallas kernel that fuses the up projection, GELU, and down
projection:
\[
  h_{\mathrm{tile}} = x W_{\mathrm{up},\mathrm{tile}},\qquad
  h_{\mathrm{tile}} = \mathrm{gelu}(h_{\mathrm{tile}}),\qquad
  \mathrm{acc} \mathrel{+}= h_{\mathrm{tile}} W_{\mathrm{down},\mathrm{tile}}.
\]
The fused kernel streams over tiles of \(F\), so the full \([M,F]\) intermediate
is never written to and reread from global memory.

The best-performing version was a per-expert Pallas backend. Instead of passing
all expert weights into a single custom call, it launches one call per expert
and passes only that expert's
\[
  W_{\mathrm{up},e}\in\mathbb{R}^{D\times F},\qquad
  W_{\mathrm{down},e}\in\mathbb{R}^{F\times D}.
\]
This reduced VMEM pressure compared with the compact all-expert kernel.

The strongest measured configuration was:
\[
  D = 512,\qquad \mathrm{mlp\_ratio}=8,\qquad F=4096,\qquad
  E=4,\qquad M=6656.
\]
The benchmark used batch size \(B=512\), sequence length \(T=13\), top-1
routing, and float32. The winning tile sizes were:
\[
  \mathrm{BLOCK\_M}=256,\qquad \mathrm{BLOCK\_F}=128.
\]

\begin{table}[h]
\centering
\begin{tabular}{lrrrr}
\toprule
Configuration & Baseline (ms) & Fused (ms) & Speedup & Max Abs. Diff. \\
\midrule
\(D=384,\ r=4\), compact Pallas & 3.52 & 3.82 & \(0.92\times\) & 0.0 \\
\(D=384,\ r=8\), compact Pallas & 5.11 & 4.42 & \(1.15\times\) & 0.0 \\
\(D=512,\ r=4\), compact Pallas & 5.16 & 5.26 & \(0.98\times\) & 0.0 \\
\(D=512,\ r=8\), per-expert Pallas & 11.97 & 7.76 & \(1.54\times\) & 0.0 \\
\bottomrule
\end{tabular}
\caption{Full forward-pass timings for baseline \(\mathrm{ragged\_dot}\) and
fused Pallas MoE FFN implementations. Here \(r=\mathrm{mlp\_ratio}\).}
\end{table}

We also isolated the MoE FFN itself for the winning shape
\((M,D,F,E)=(6656,512,4096,4)\):

\begin{table}[h]
\centering
\begin{tabular}{lrrrr}
\toprule
Routing Pattern & Baseline (ms) & Pallas (ms) & Speedup & Max Abs. Diff. \\
\midrule
Balanced \([1664,1664,1664,1664]\) & 1.572 & 0.993 & \(1.58\times\) & 0.0 \\
Skewed \([512,1024,2048,3072]\) & 1.560 & 0.966 & \(1.61\times\) & 0.0 \\
\bottomrule
\end{tabular}
\caption{Isolated MoE FFN timings for the winning Pallas kernel shape.}
\end{table}

The speedup appears when the hidden activation is large enough that avoiding
its materialization dominates custom-kernel overhead. For \(D=512\) and
\(r=8\), the baseline hidden activation has shape
\[
  [M,F] = [6656,4096],
\]
or
\[
  6656\cdot 4096 = 27{,}262{,}976
\]
float32 values, roughly \(109\) MB. The baseline path writes the up-projection
output, reads it for GELU, writes the GELU result, and reads it again for the
down projection. The fused kernel instead keeps only a tile such as
\([256,128]\) live inside the kernel and consumes it immediately.

At smaller shapes, especially \(D=384,r=4\), the baseline
\texttt{jax.lax.ragged\_dot} remains competitive. The saved intermediate is not
large enough to overcome routing, packing, custom-call, and scatter/gather
overheads. For wider shapes, fusion should become more valuable, but the current
implementations hit TPU VMEM limits if too much expert weight data is staged at
once. Output-dimension tiling can make larger cases compile, but it recomputes
the up projection and GELU for each output tile, which can erase the benefit.

The practical conclusion is that the fused Pallas MoE FFN is worthwhile once
\(F\) is large relative to \(D\). In our experiments, the clearest win was
\(D=512,r=8,F=4096\), where the per-expert Pallas kernel achieved about
\(1.54\times\) end-to-end forward speedup and about \(1.6\times\) isolated FFN
speedup with exact float32 output parity.

\end{document}
